\subsection{A seven-cell tree obstruction}\label{u:section}

We prove a slightly sharper upper bound than required: every controller
with at most two states has a trap of size at most seven that is either
an induced path or an induced tree with one vertex of degree three.
Throughout this section, an action is written
$\delta(q,M)=(d,r)$, with direction first. We call a controller a
\emph{path survivor} if it succeeds from every pair of state-zero starts
on every induced grid path with at most seven cells.

Two elementary observations will certify the traps. If two executions
stay in regions whose cross-distances are at least two, they cannot
succeed. If instead their regions are merely disjoint, the same conclusion
holds when their starts have equal checkerboard color: compulsory
synchronous motion preserves equality of their colors, whereas adjacent
cells have different colors. In every application below the regions
contain the initial positions, or an explicit entry prefix is given.

All compass normalizations act on the controller as well as on the
maze. For $g\in D_4$, define
\[
 (g\delta)(q,gM)=(gd,r)\quad\text{when }\delta(q,M)=(d,r).
\]
Induction on time conjugates the complete tagged executions under
$(x,q)\mapsto(gx,q)$. Manhattan distance and initial state zero are
preserved. In particular, a trap for the transported controller gives
a trap for the original controller by applying $g^{-1}$.

\subsubsection{Straight-axis rigidity}

\begin{lemma}[Axis reduction]\label{u:axis}
For a path survivor, one simultaneous compass symmetry makes the
straight rows
\[
\begin{array}{ll}
\delta(0,\mathrm{NS})=(\Sdir,f_v(0)),&
\delta(1,\mathrm{NS})=(\N,f_v(1)),\\
\delta(0,\mathrm{EW})=(\W,f_h(0)),&
\delta(1,\mathrm{EW})=(\E,f_h(1)).
\end{array}
\]
Each map $f=f_v,f_h$ is a constant map $C_0,C_1$ or the identity $I$.
Write $L_q$ for the next-state bit at the negative endpoint of an axis
and $R_q$ for that at its positive endpoint. The necessary restrictions are
\[
\begin{array}{c|l}
f&\text{endpoint restrictions}\\ \hline
C_0&L_0=1,\ R_0=0,\\
I&L_0=1,\ (R_0,R_1)\ne(1,1),\\
C_1&L_1=1,\ (R_0,R_1)\ne(1,1).
\end{array}\tag{\ref{u:axis}a}
\]
All unlisted endpoint bits remain arbitrary.
\end{lemma}

\begin{proof}
For perpendicular unit directions $d,e$, the seven cells
\[
 U(d,e)=\{0,e,2e,d,2d,2e+d,2e+2d\}
\]
form an induced path with two parallel two-edge arms pointing in direction
$d$. If both straight rows on an axis choose $d$, starts at $d$ and
$2e+d$ stay on translated terminal edges, regardless of all next-state
bits. Their distance is always two. Thus the straight directions are
opposite, and independent coordinate reflections give the displayed
normalization without exchanging the states.

Index a straight path by increasing coordinates. If $f=(1,0)$, the
six-cell path $0,\ldots,5$ with starts $2,4$ gives the two tagged cycles
$(2_0,1_1)$ and $(4_0,3_1)$; the agents remain two cells apart.
If $f=C_0$ and $R_0=1$, the five-cell path $0,\ldots,4$ with starts
$0,4$ confines the right agent to $(4_0,3_1)$ and the left to $0,1,2$.
Indeed the left word is $(0_0,1_0)^\infty$ when $L_0=0$, and
$(0_0,1_1,2_0,1_0)^\infty$ when $L_0=1$.
The regions are disjoint and the starts have equal color.
If $f=C_1$ and $L_1=0$, use the same line with starts $1,3$.
The left cycle is $(1_0,0_1)$; the right prefix
$3_0\to2_1\to3_1$ enters the closed region $2,3,4$.
At its positive endpoint, either return bit keeps it in that region.
Again equal colors exclude adjacency.

Two terminal-edge tests finish the restrictions. If $L_{f(0)}=0$,
an inner state-zero start has cycle
$I_0\to E_{f(0)}\to I_0$ at a negative terminal edge. Two copies
in $U(d,e)$, with $d$ negative and inner starts, give a trap.
If $R_0=R_{f(1)}=1$, an endpoint start at a positive terminal edge has
word $E_0\to I_1\to E_{f(1)}\to I_1\to\cdots$.
Clone these words in the positive-pointing $U(d,e)$, with endpoint
starts. Both constructions keep the agents two units apart.
The four possible two-bit maps, together with these exclusions, give
exactly the necessary restrictions displayed above.
\end{proof}

For singleton masks write
$\delta(q,\N)=(\N,n_q)$, and similarly $e_q,s_q,w_q$.
Thus $(L,R)=(n,s)$ vertically and $(e,w)$ horizontally.
For perpendicular directions $u,v$, put
\[
 L(u,v;a,b)=\{C,u_1,\ldots,u_a,v_1,\ldots,v_b\},
 \qquad C=(0,0),\quad d_j=jd.
\]
This is an induced path with $a+b+1$ cells.

\begin{lemma}[Reset pocket]\label{u:pocket}
Suppose a straight state map resets to $\kappa$, and let $d$ be the
direction selected by state $\kappa$. On a terminal arm
$C,d_1,d_2,d_3$, a state-zero start at $d_2$ never reaches $C$.
\end{lemma}

\begin{proof}
The tagged set
$\{(d_1)_\kappa,(d_2)_0,(d_2)_1,(d_3)_\kappa\}$ is closed.
At $d_2$, either action reaches $d_1$ or $d_3$ in state $\kappa$.
From $(d_1)_\kappa$ the next move goes outward to $(d_2)_\kappa$;
the endpoint $(d_3)_\kappa$ returns to $d_2$ in either state.
The set contains the initial tagged position.
\end{proof}

% Exact geometry: C=(0,0), S_j=(0,-j); other neighbours of C omitted.
\begin{figure}[tbp]
\centering
\begin{tikzpicture}[maze,x=8mm,y=8mm]
 \fill[rvGray!60,rounded corners=2mm] (-.28,-3.27) rectangle (.28,-.72);
 \draw[quiet] (0,0)--(0,-1);
 \draw[supportA] (0,-1)--(0,-2)--(0,-3);
 \node[gate,label=above:{$C=(0,0)$}] at (0,0) {};
 \node[cell,label=left:{$S_1$}] at (0,-1) {};
 \node[startA,label=left:{$(S_2)_0$}] at (0,-2) {};
 \node[cell,label=left:{$S_3$}] at (0,-3) {};
 \node[paneltitle] at (1.35,-.25) {Closed reset pocket};
 \node[note] at (1.35,-1.02) {$f=C_0,\quad d=\Sdir$};
 \node[note] at (1.35,-1.82) {$\mathcal P=\{(S_1)_0,(S_2)_0,(S_2)_1,(S_3)_0\}$};
 \node[note] at (1.35,-2.65) {Start in $\mathcal P$; the corner is never visited.};
 \mazeCompass{9.2}{-2.1}
\end{tikzpicture}
\caption{The closed tagged pocket for reset state $0$ and direction $\Sdir$; unshown neighbours of $C$ do not affect it. Other reset directions transport both the geometry and the controller directions together.}
\label{f:reset-pocket}
\end{figure}


\begin{lemma}[Straight-axis rigidity]\label{u:rigidity}
For a path survivor in the frame of Lemma~\ref{u:axis}, both
straight state maps are the identity. Consequently
\[
 n_0=e_0=1,
 \qquad (s_0,s_1),(w_0,w_1)\ne(1,1).
\]
\end{lemma}

\begin{proof}
If both axes reset, choose their outward reset directions $d_v,d_h$.
On $L(d_v,d_h;3,3)$ start at $(d_v)_2,(d_h)_2$.
Lemma~\ref{u:pocket} confines the agents to different three-cell
arms, whose cross-distances are at least two. This treats all four
choices of the two reset bits.

If exactly one axis resets, interchange the axes simultaneously to
make it vertical. The two remaining possibilities are
$C_0/I$ and $C_1/I$. Lemmas~\ref{u:mixed0} and~\ref{u:mixed1}
give explicit path obstructions for each in
Appendix~\ref{u:appendix-mixed}. Their proofs use reset pockets,
closed arms, and short transition words; all endpoint bits allowed by
Lemma~\ref{u:axis} are retained. Thus neither mixed possibility can
survive. The endpoint conclusions now follow from Lemma~\ref{u:axis}.
\end{proof}

\subsubsection{Corner compatibility and a nine-entry pattern}

\begin{lemma}[Corner arrival]\label{u:arrival}
Use the normalized straight directions of Lemma~\ref{u:axis}.
Let $C,I,E$ be a terminal arm of two edges, with $I$ a straight cell
and direction $d$ from $C$ into the arm. Every arrival at $C$ from
this arm has state $a_d=f(q_{-d})$, where state $q_{-d}$ selects the
straight direction toward $C$. If the corner action in state $a_d$
points back into the arm, a state-zero start at $I$ remains in
$C,I,E$. For identity straight maps,
\[
a_\N=a_\E=0,\qquad a_\Sdir=a_\W=1.
\]
\end{lemma}

\begin{proof}
The last move $I\to C$ gives the arrival formula. Endpoints return
to $I$, straight moves stay in the arm or reach $C$ in that state,
and the specified corner row sends such an arrival back to $I$.
This proves closure from the stated initial position, with no
restriction on endpoint or corner output bits.
\end{proof}

Call the condition in this lemma a \emph{locked arm}. Two corners
two units apart, joined by their middle cell, can carry perpendicular
two-edge arms on either side of the connector. If both arms are locked,
starts at their inner cells stay in parallel rows or columns two
units apart. Each construction is an induced seven-cell path.

\begin{lemma}[Corner compatibility]\label{u:corners}
After one further diagonal reflection, a path survivor satisfies
\[
\begin{array}{ll}
\delta(0,\mathrm{NE})=(\E,1),&
\delta(0,\mathrm{ES})=(\Sdir,0),\\
\delta(1,\mathrm{ES})=(\E,1),&
\delta(1,\mathrm{SW})=(\Sdir,0).
\end{array}\tag{\ref{u:corners}a}
\]
\end{lemma}

\begin{proof}
The diagonal reflection exchanges the axes, preserves the four
identity straight rows, and lets us choose the direction of
$\delta(0,\mathrm{NE})$ to be $\E$. Its east arm is locked.
The following three paired-arm obstructions force the other directions.
First put an NE corner at $(0,0)$ with east arm and an SW corner at
$(0,2)$ with west arm, joined through $(0,1)$. If SW in state 1
chose west, both arms would be locked. Hence SW in state 1 chooses
south. Next replace the upper corner by ES with an east arm; if ES
in state 0 chose east, both east arms would be locked. Thus ES in
state 0 chooses south. Finally use $U(\Sdir,\E)$: its two corners
have masks ES and SW. The SW south arm is locked, so ES in state 1
cannot also choose south. It must choose east. Starts in all three
constructions are the inner cells of the two arms.

These direction conclusions leave four output bits. The four
six-cell obstructions in Lemma~\ref{u:bit-laws} force them to be,
respectively, $1,0,1,0$. Each argument uses disjoint regions and
equal-color starts, and none assumes any of the other three bit
conclusions.
\end{proof}

% Exact six-cell grid geometry; highlighted regions are closed, not fixed cycles.
\begin{figure}[tbp]
\centering
\begin{tikzpicture}[maze,x=9mm,y=9mm]
 \coordinate (P) at (1,2); \coordinate (U) at (0,2);
 \coordinate (M) at (0,1); \coordinate (C) at (0,0);
 \coordinate (H) at (1,0); \coordinate (Eend) at (2,0);
 \draw[quiet] (M)--(C);
 \draw[supportB] (P)--(U)--(M);
 \draw[supportA] (C)--(H)--(Eend);
 \node[cell,label=right:{$P$}] at (P) {};
 \node[startB,label=above left:{$U_0$}] at (U) {};
 \node[cell,label=left:{$M_1$}] at (M) {};
 \node[startA,label=below left:{$C_0$}] at (C) {};
 \node[cell,label=below:{$H$}] at (H) {};
 \node[cell,label=right:{$E$}] at (Eend) {};
 \node[paneltitle] at (3.2,2.0) {Two closed regions};
 \node[note] at (3.2,1.25) {Forbidden row: $\delta(0,ES)=(S,1)$};
 \node[note] at (3.2,.55) {$C_0:\ \{C,H,E\}\qquad U_0:\ \{U,M,P\}$};
 \node[note] at (3.2,-.15) {Disjoint visits; equal checkerboard colour.};
 \mazeCompass{10.1}{1.15}
\end{tikzpicture}
\caption{Corner obstruction under the stated forbidden row and preceding rigidity hypotheses. The $U$-started execution reaches $M$ only in state $1$; disjointness excludes coincidence and parity excludes adjacency.}
\label{f:corner-path}
\end{figure}


\begin{lemma}[Nine-entry pattern]\label{u:nine}
Every path survivor admits a simultaneous compass frame in which
the following entries hold:
\[
\begin{array}{c|c|c@{\qquad}c|c|c}
M&q&\delta(q,M)&M&q&\delta(q,M)\\ \hline
\N&0&(\N,1)&\mathrm{NE}&1&(\N,1)\\
\E&0&(\E,1)&\mathrm{ES}&0&(\Sdir,0)\\
\W&1&(\W,1)&\mathrm{ES}&1&(\E,1)\\
\mathrm{EW}&1&(\E,1)&\mathrm{SW}&1&(\Sdir,0)\\
\mathrm{NE}&0&(\E,1)&&&
\end{array}\tag{\ref{u:nine}a}
\]
No junction entry is prescribed.
\end{lemma}

\begin{proof}
Only $w_1=1$ and $\delta(1,\mathrm{NE})=(\N,1)$ remain to be
proved. Use the path, in endpoint order,
\[
 (1,0),(0,0),(0,1),(0,2),(1,2),(2,2),(2,1),
\]
with state-zero starts $A=(0,0)$ and $Q=(2,1)$.
Put $P=(1,0)$, $M=(0,1)$, and $D=(2,2)$.
The right agent has cycle $Q_0\to D_1\to Q_0$.
The left starts $A_0\to P_1$. If $w_1=0$, it returns to $A_0$
and repeats. If $w_1=1$, it reaches $A_1$. Should the NE state-one
row choose east, $A,P$ is closed; should it be $(\N,0)$, its only
extra excursion is $A_1\to M_0\to A_0$. In every bad case the
left stays in $A,P,M$, all at distance at least two from $Q,D$.
Thus the two remaining entries are forced, and the other seven
come from Lemmas~\ref{u:rigidity} and~\ref{u:corners}.
\end{proof}

% Exact seven-cell path. Left highlight is a conditional closed region.
\begin{figure}[tbp]
\centering
\begin{tikzpicture}[maze,x=9mm,y=9mm]
 \coordinate (P) at (1,0); \coordinate (A) at (0,0);
 \coordinate (M) at (0,1); \coordinate (T) at (0,2);
 \coordinate (H) at (1,2); \coordinate (D) at (2,2);
 \coordinate (Q) at (2,1);
 \draw[quiet] (M)--(T)--(H)--(D);
 \draw[supportA] (P)--(A)--(M); \draw[supportB] (Q)--(D);
 \node[cell,label=below:{$P$}] at (P) {};
 \node[startA,label=below left:{$A_0$}] at (A) {};
 \node[cell,label=left:{$M$}] at (M) {};
 \node[cell] at (T) {}; \node[cell] at (H) {};
 \node[cell,label=above:{$D_1$}] at (D) {};
 \node[startB,label=right:{$Q_0$}] at (Q) {};
 \node[paneltitle] at (3.2,2.0) {The final two entries};
 \node[note] at (3.2,1.25) {$Q_0\longleftrightarrow D_1$ is forced.};
 \node[note] at (3.2,.55) {A failed entry confines $A_0$ to $\{A,P,M\}$.};
 \node[note] at (3.2,-.15) {$\operatorname{dist}_1(\{A,P,M\},\{D,Q\})\ge2$};
 \mazeCompass{10.1}{1.15}
\end{tikzpicture}
\caption{The final seven-cell path, anchored at $A=(0,0)$. If either remaining W9 entry fails, the left execution stays separated from the forced right-hand tagged two-cycle, including the transient.}
\label{f:final-path}
\end{figure}


\subsubsection{Completion at a junction}

The last step uses only the nine entries of Lemma~\ref{u:nine}; all
other rows remain arbitrary.

\begin{lemma}[Junction completion]\label{u:junction}
Every legal completion of the nine-entry pattern has a trap on one
of the two seven-cell induced trees $G_1,G_2$ below.
\end{lemma}

\begin{proof}
Write $\delta(1,\E)=(\E,\epsilon)$ and
$u=\delta(0,\mathrm{NSW})$, $v=\delta(1,\mathrm{NSW})$.
The common cells of the two trees are
\[
A=(0,1),\quad B=(1,0),\quad J=(1,1),\quad
U=(1,2),\quad V=(2,2).
\]
For $G_1$ add $D=(2,0),F=(3,0)$ and start at $A,B$.
For $G_2$ add $Q=(3,1),R=(3,2)$ and start at $B,Q$.
All starts have state zero. The edge sets are
\[
\begin{array}{c|l}
G_1&JA,JB,JU,UV,BD,DF,\\
G_2&JA,JB,JU,UV,VR,RQ.
\end{array}
\]
Direct coordinate comparison gives no other unit-distance pair.
Each tree has one degree-three vertex $J$, with mask NSW.
Choose $G_2$ precisely when
\[
 v=(\Sdir,0)\quad\text{or}\quad
 \bigl[u=(\Sdir,0)\text{ and }
 v\in\{(\N,0),(\W,1)\}\bigr].\tag{\ref{u:junction}a}
\]
Choose $G_1$ otherwise.

On $G_1$, the B-started agent follows
$B_0\to D_1\to F_1\to D_1\to\cdots$.
The A-started agent begins $A_0\to J_1$. Its relevant nonjunction
transitions are
\[
\begin{gathered}
A_0\to J_1,\quad A_1\to J_\epsilon,\quad B_1\to J_1,\\
U_0\to J_0,\quad U_1\to V_1,\quad V_1\to U_1.
\end{gathered}
\]
If neither $u$ nor $v$ is $(\Sdir,0)$, these transitions and the
allowed junction departures close the tagged set
\[
\{A_0,A_1,J_0,J_1,B_1,U_0,U_1,V_1\}.
\]
The remaining $G_1$ cases have $u=(\Sdir,0)$ and
$v\in\{(\N,1),(\Sdir,1),(\W,0)\}$.
They send the A-started agent, respectively, to the edge $UV$, to
the state-one cycle on $JB$, or to $J_1\leftrightarrow A_0$.
It never queries $J_0$, so the row $u$ is harmless.
Thus from time one the two agents lie on disjoint regions, one on
$DF$ and the other outside it. At time zero they are distinct and
have equal checkerboard color, which excludes adjacency at every
time. Hence $G_1$ is a trap in every case where it was selected.

On $G_2$ the Q-started cycle is always $Q_0\leftrightarrow R_1$.
The selection condition gives exactly the following B-started words:
\[
\begin{array}{l|l}
v=(\Sdir,0)&(B_0,J_1)^\infty\\
u=(\Sdir,0),\ v=(\N,0)&(B_0,J_1,U_0,J_0)^\infty\\
u=(\Sdir,0),\ v=(\W,1),\ \epsilon=0
 &(B_0,J_1,A_1,J_0)^\infty\\
u=(\Sdir,0),\ v=(\W,1),\ \epsilon=1
 &B_0,(J_1,A_1)^\infty.
\end{array}
\]
Every indicated move is fixed by the nine entries or the displayed
choice of $u,v,\epsilon$. At each integer time the distance from
the other agent is exactly three, as follows directly from the
coordinates. This includes all transients and finishes the proof.
\end{proof}

% Both embeddings retain the exact coordinates from the original junction proof.
\begin{figure}[tbp]
\centering
\begin{tikzpicture}[maze,x=8.5mm,y=8.5mm]
 \begin{scope}
  \coordinate (A) at (0,1); \coordinate (B) at (1,0); \coordinate (J) at (1,1);
  \coordinate (U) at (1,2); \coordinate (V) at (2,2); \coordinate (D) at (2,0); \coordinate (F) at (3,0);
  \draw[edge] (A)--(J)--(B)--(D)--(F); \draw[edge] (J)--(U)--(V);
  \draw[supportB] (D)--(F); \draw[transientB] (1.15,-.15)--(1.83,-.15);
  \node[startA,label=left:{$A_0$}] at (A) {};
  \node[startB,label=below:{$B_0$}] at (B) {};
  \node[cell,label=left:{$J$}] at (J) {}; \node[cell,label=above:{$U$}] at (U) {};
  \node[gate,label=above:{$V$}] at (V) {}; \node[cell,label=below:{$D_1$}] at (D) {};
  \node[cell,label=right:{$F_1$}] at (F) {};
  \node[paneltitle] at (-.35,2.95) {(a) $G_1$: arms $1,2,3$};
  \node[note] at (-.35,-.92) {$B_0\to D_1\leftrightarrow F_1$; $M(V)=\W$.};
 \end{scope}
 \draw[rvGray,line width=.65pt] (4.45,-1.13)--(4.45,3.1);
 \begin{scope}[xshift=48mm]
  \coordinate (A) at (0,1); \coordinate (B) at (1,0); \coordinate (J) at (1,1);
  \coordinate (U) at (1,2); \coordinate (V) at (2,2); \coordinate (Q) at (3,1); \coordinate (R) at (3,2);
  \draw[edge] (A)--(J)--(B); \draw[edge] (J)--(U)--(V)--(R);
  \draw[supportB] (Q)--(R);
  \node[cell,label=left:{$A$}] at (A) {}; \node[startA,label=below:{$B_0$}] at (B) {};
  \node[cell,label=left:{$J$}] at (J) {}; \node[cell,label=above:{$U$}] at (U) {};
  \node[gate,label=above:{$V$}] at (V) {}; \node[startB,label=right:{$Q_0$}] at (Q) {};
  \node[cell,label=right:{$R_1$}] at (R) {};
  \node[paneltitle] at (-.35,2.95) {(b) $G_2$: arms $1,1,4$};
  \node[note] at (-.35,-.92) {$Q_0\leftrightarrow R_1$; $M(V)=\E\W$.};
 \end{scope}
 \mazeCompass{10.6}{1.3}
\end{tikzpicture}
\caption{Junction completions with $J=(1,1)$ and $M(J)=\N\Sdir\W$. The highlighted right-hand shuttle is fixed in each selected case; the other execution is governed by the exact alternatives in the proof.}
\label{f:junctions}
\end{figure}


\begin{proposition}[Upper bound in a narrow tree class]\label{u:upper}
Every controller with at most two states has a trap with at most
seven cells that is either an induced path or an induced tree with
exactly one vertex of degree three and no larger degree. In particular,
$F(2)\le7$ and $F_{\mathrm{grid\ trees}}(2)\le7$.
\end{proposition}

\begin{proof}
A two-state controller either already has a path trap of size at
most seven, or is a path survivor. In the second case apply
Lemma~\ref{u:nine}, then Lemma~\ref{u:junction}, and transport the
resulting tree and starts back by the inverse compass symmetry.
A one-state controller can be extended by a legal unused second
state without changing any execution from its distinguished initial
state. The two-state argument therefore covers it as well.
\end{proof}
